% !PRINTED 231
\weberpage{231}
%<<CONT
\weq{12}{(-1)^{\frac{n^2-1}{8}}\,2^{\frac{n^2(n-1)}{3}}P_{11} = \sqrt{\varkappa\varkappa'}^{\frac{n-1}{3}}\mathop{\Pi}\limits_{1,\frac{n-1}{2}}^{h}\frac{\operatorname{sn}h\Omega\,D(\operatorname{sn}^2h\Omega)^2}{(\operatorname{cn}h\Omega)^{\frac{2n^2+1}{3}}(\operatorname{dn}h\Omega)^{\frac{2n^2+1}{3}}}.}

Soll als Rationalitätsbereich der der rationalen Zahlen und rationalen Funktionen von $\varkappa^2$ aufrecht erhalten werden, so schreibt man die letzte Gleichung besser in der Form [vgl. § 54, (4)]:
\weq{13}{\begin{aligned}
P_{11}\gamma_2(\omega)^{\frac{n-1}{2}} &= (-1)^{\frac{n^2-1}{8}}\,2^{-\frac{(n^2-4)(n-1)}{3}}(\varkappa\varkappa')^{-\frac{n-1}{2}}(1-\varkappa^2\varkappa'^2)^{\frac{n-1}{2}}\\
&\quad\times \mathop{\Pi}\limits_{1,\frac{n-1}{2}}^{h}\frac{\operatorname{sn}h\Omega\,D(\operatorname{sn}^2h\Omega)^2}{(\operatorname{cn}h\Omega)^{\frac{2n^2+1}{3}}(\operatorname{dn}h\Omega)^{\frac{2n^2+1}{3}}},
\end{aligned}}
und schließt daraus auf den folgenden Satz:

\gesperrt{Wenn $n$ nicht durch 3 teilbar ist, so sind $P_{00}$, $P_{01}$, $P_{10}$, $P_{11}^2\gamma_2(\omega)^{n-1}$, und wenn $n$ ein Quadrat ist, auch $P_{11}\gamma_2(\omega)^{\frac{n-1}{2}}$ Wurzeln von Transformationsgleichungen.}

\begin{center}
\textbf{§ 68.} \textbf{\gesperrt{Zweite Darstellung der Wurzeln der Transformationsgleichungen.}}
\end{center}

Wenn wir zunächst eine der Reihen ins Auge fassen, nämlich die, zu der die Wurzel $x_{1,0}$ gehört, also $\mu = 1$, $\mu' = 0$ setzen, so können wir auf (2), (3) des vorigen Paragraphen die Formeln (20), (21), (10), § 34 anwenden. Setzt man dort, wenn $\nu$ ungerade ist, $n - \nu$ an Stelle von $\nu$ und benutzt die Formel
\[
\vartheta_{g_1g_2}\left(\frac{n-\nu}{n}\right) = (-1)^{g_1(g_2+1)}\vartheta_{g_1g_2}\left(\frac{\nu}{n}\right),
\]
dann kommen in diesen Formeln nur die geraden $\nu$ vor, und wenn man also $\nu = 2h$ setzt, so kann man die dortigen Formeln (20), (21) auch so schreiben:
\[
\begin{aligned}
\sqrt{n}\,\eta(n\omega)\eta(\omega)^{\frac{n-3}{2}} &= \mathop{\Pi}\limits_{1,\frac{n-1}{2}}^{h}\vartheta_{11}\!\left(\frac{2h}{n}\right),\\
f(n\omega)\eta(\omega)^{\frac{n-1}{2}} &= f(\omega)\mathop{\Pi}\limits_{1,\frac{n-1}{2}}^{h}\vartheta_{00}\!\left(\frac{2h}{n}\right),\\
f_1(n\omega)\eta(\omega)^{\frac{n-1}{2}} &= f_1(\omega)\mathop{\Pi}\limits_{1,\frac{n-1}{2}}^{h}\vartheta_{01}\!\left(\frac{2h}{n}\right),\\
f_2(n\omega)\eta(\omega)^{\frac{n-1}{2}} &= \left(\frac{2}{n}\right)f_2(\omega)\mathop{\Pi}\limits_{1,\frac{n-1}{2}}^{h}\vartheta_{10}\!\left(\frac{2h}{n}\right),
\end{aligned}
\]
%CONT>>
